\section{Facilities, Equipment, and Other Resources}

\subsection*{University of Texas at El Paso (UTEP)}
PI Rahman is an assistant professor in the Department of Computer Science at the University of Texas at El Paso (UTEP). The Department of Computer Science maintains strong research and educational programs with a particular emphasis on cybersecurity and artificial intelligence (AI). The department offers a specialization in cybersecurity as part of its M.S. in Software Engineering program, as well as dedicated Bachelor's and Master's programs in AI, underscoring its commitment to advancing research and education in these critical areas. The PI's faculty office is located in the Chemistry and Computer Science Building at UTEP.



\subsection*{T3S Lab at UTEP}
The T3S Lab, led by Co-PI Armanuzzaman and available at \url{https://t3slab.github.io/}, is housed in the Department of Computer Science at UTEP and supports collaborative research in embedded systems security, systems and software security, hardware security, FPGA security, and agentic systems security. The lab currently supports one Ph.D. student, two M.S. students, and two B.S. students. The lab has high-speed wired and wireless Internet access and is equipped with a dedicated high-performance workstation for development, benchmarking, and reproducible experimentation. In addition, the lab maintains a suite of embedded development boards and security evaluation tools, including the Renesas RA8M1 Cortex-M85 evaluation board \cite{renesasra8m1}, Cortex-M33 development boards, the ARM Versatile Express V2M-MPS3 prototyping board \cite{ARMMPS3}, a Dell PowerEdge R570 server with an NVIDIA L4 PCIe GPU (72W, 24 GB, passive, single-wide full-height), and the NewAE ChipWhisperer HuskyPlus kit \cite{chipwhispererhuskyplus}. These resources support research on microcontroller security, side-channel analysis, firmware attestation, hardware-enabled trust, FPGA security, and secure agentic system design. The lab also includes standard desktop workstations and software development tools for system prototyping and evaluation.

\subsection*{High Performance Computing (HPC) Resources}
UTEP maintains centralized HPC resources for large-scale data analysis, scientific computing, and AI/ML workloads, with support from professional staff and integration with the Texas Advanced Computing Center (TACC) at UT Austin. The campus operates three clusters: JAKAR (2,880 cores, 28 TB RAM), PARO (913 cores, 11.5 TB RAM), and PUNAKHA, UTEP's GPU-accelerated AI cluster. PUNAKHA consists of two NVIDIA DGX/HGX nodes with twelve NVIDIA H100 GPUs, 3 TB system memory, 30 TB local storage, and NVIDIA NDR200/200GbE InfiniBand connectivity. It runs Red Hat Enterprise Linux 9.2 with OpenHPC 3, uses Docker and SLURM for workload management, and supports Multi-Instance GPU (MIG) partitioning for flexible GPU allocation. These systems provide the compute capacity needed for large-scale security and AI research.

\subsection*{Institutional Designations}
UTEP is classified as an R1 research institute under the Carnegie Classification of Institutions of Higher Education, reflecting its very high research activity. The university is designated by the National Security Agency (NSA) and the Department of Homeland Security (DHS) as a National Center of Academic Excellence in Cyber Defense (CAE-CD). In 2016, UTEP was also designated as a National Center of Academic Excellence in Cyber Operations (CAE-CO). These designations underscore the institution's strength in cybersecurity research and education and provide a strong foundation for research initiatives in AI and security.

\subsection*{Unfunded Collaborations}
The project benefits from unfunded industrial and academic collaborations that provide expertise and complementary resources.

\textbf{Industrial Collaborations:}
\begin{itemize}
    \item \textbf{Peter Linder, Cofounder and Chief Product Architect, Octavo Systems}
    \begin{itemize}
        \item Contributes expertise in embedded systems design and development, particularly for FPGA-based SoC platforms.
        \item Provides perspective on practical industry workflows, hardware/software co-design, and secure implementation strategies.
    \end{itemize}

    \item \textbf{Malav Vyas, Security Researcher, Palo Alto Networks}
    \begin{itemize}
        \item Contributes expertise in SCADA and industrial control system security.
        \item Brings industry insight on IoT security, digital forensics, and applied security analysis in real-world systems.
    \end{itemize}
\end{itemize}
